% GENERATED FILE. Edit canonical lessons, then run npm run generate:books.
% canonical-source-hash: fnv1a64:6981c0a567ac9e97
% canonical-lessons: ZH-C95-weishenme, ZH-W95-wei, ZH-C95-youshi, ZH-C95-duche, ZH-C95-wang, ZH-C95-practice

\chapter{Why?}
\label{ch:zh-why}
\hlchaptermodality{95}

\begin{chapteropening}
\textbf{By the end of this chapter:} \emph{I can ask why.}

Asking why with \zh{为什么}, and three everyday answers to it, being busy, the traffic and forgetting, built on one new written character: \zh{为}.
\end{chapteropening}

\section[\texorpdfstring{\zh{为什么} (wèishénme)}{为什么 (wèishénme)}]{\zh{为什么} (wèishénme) — why?}

\label{lesson:ZH-C95-weishenme}



\begin{quote}
A new chapter, and it opens with the question that wants a reason.

\begin{itemize}
\raggedright
  \item \emph{Recall:} say \zh{我们都是学生} — \emph{wǒmen dōu shì xuésheng}
  \item \emph{Recall:} say \emph{tā hěn shòu}, he's thin — \textbf{R1}, the chapter you have finished
  \item \emph{Recall:} say \emph{chídào}, to be late — \textbf{R2}
\end{itemize}
\end{quote}

\subsection*{You'll want to know: \zh{为什么}}
\textbf{\zh{为什么}} — \emph{wèishénme} — \textquotedblleft{}why?\textquotedblright{}.

\begin{quote}
\textbf{\zh{你为什么不去}} \emph{nǐ wèishénme bú qù} — Why aren't you going?
\end{quote}

The next lesson puts the new character in your hand.

\begin{cousinweb}
\textbf{\zh{为什么}} (\emph{wèishénme}) is built from \zh{为} 'for' and \zh{什么} 'what': for what? \zh{为} is inherited from Middle Chinese \emph{hjweH} (Guangyun, in Baxter's transcription), glossed 'to help'; 'for' is the sense of acting on somebody's behalf, which the gloss does not spell out. Where \zh{什么} comes from is not settled.
\end{cousinweb}

\begin{grammarlens}[title={What you've built}]
\zh{为什么} sits after the person and before the verb: \zh{你}, \zh{为什么}, \zh{不去}. \zh{不} before the fall of \zh{去} rises: \emph{bú qù}.
\end{grammarlens}

\subsection*{Your turn}
\begin{itemize}
\raggedright
  \item \emph{Say it:} \emph{wèishénme}
  \item \emph{Say it:} \emph{wèishénme}, once more
  \item \emph{Say it:} \zh{你为什么不去} — \emph{nǐ wèishénme bú qù} — Why aren't you going?
  \item \emph{Say it:} \emph{nǐ wèishénme chídào} — Why are you late?
  \item \emph{Recall:} say \emph{chídào}, then ask why
\end{itemize}

\begin{tcolorbox}[breakable,colback=teal!4,colframe=teal!35!black,title={Before you move on}]
What does \zh{为什么} ask? (\textquotedblleft{}Why?\textquotedblright{}.) Where does it go? (\textbf{After the person, before the verb}.)
\end{tcolorbox}
\section[\texorpdfstring{\zh{为} (wèi)}{为 (wèi)}]{\zh{为} — a dot, a long fall, the hook of \zh{男}, a dot}

\label{lesson:ZH-W95-wei}



\begin{quote}
The word first, then the character inside it.

\begin{itemize}
\raggedright
  \item \emph{Recall:} say \emph{wèishénme} — \textbf{R1}, one lesson back, before it decays
\end{itemize}

\begin{quote}
\zh{男} \zh{点}
\end{quote}

\emph{Read it:} each one aloud, with nothing said first — \emph{nán}, \emph{diǎn}
\end{quote}

\subsection*{Script — a dot, a long fall, the hook of \zh{男}, a dot}
\hlblockfigure{\detokenize{figures/ZH-W95-wei-filmstrip.pdf}}{How \zh{为} is written, stroke by stroke}

\begin{quote}
\zh{为}
\end{quote}

\textbf{Four strokes, and the long one in the middle is a stroke your hand knows from the foot of \zh{男}.}

\textbf{A dot} at the upper left.

\textbf{A long left-falling stroke} from the top, down to the lower left.

\textbf{The hooked turn}: across, down the right side, hooking to the left, as at the foot of \zh{男}.

\textbf{A dot inside}, last.

Four strokes.

\subsection*{Your turn}
\begin{itemize}
\raggedright
  \item \emph{Write it:} \zh{为}, all four, with the model in front of you
  \item \emph{Cover:} write it again from memory, and count the pieces, not the strokes
\end{itemize}

\begin{tcolorbox}[breakable,colback=teal!4,colframe=teal!35!black,title={Before you move on}]
How many strokes in \zh{为}? (\textbf{Four}.) Where have you made the hooked turn before? (\textbf{At the foot of \zh{男}}.) Which stroke is last? (\textbf{The dot inside}.)

Source: \href{https://raw.githubusercontent.com/chanind/hanzi-writer-data/68d10a4b21150cae5e1ebbd223eed289cf32d90c/data/\%E4\%B8\%BA.json}{Hanzi Writer Data \zh{为}, pinned stroke paths and medians}.
\end{tcolorbox}
\section[yǒushì]{yǒushì — to have something on}

\label{lesson:ZH-C95-youshi}



\begin{quote}
Before the new one: read the last lesson's character, then two words for a full day.

\begin{itemize}
\raggedright
  \item \emph{Read it:} \zh{为什么}, aloud — \emph{wèishénme}
  \item \emph{Recall:} say \emph{máng}, busy — \textbf{R4}, recognition at distance
  \item \emph{Recall:} say \zh{今天} — \emph{jīntiān}, today
\end{itemize}
\end{quote}

\subsection*{You'll want to know: yǒushì}
\emph{yǒushì} — \textquotedblleft{}to have something on; to be busy\textquotedblright{}.

\begin{quote}
\emph{wǒ jīntiān yǒushì} — I've got something on today.
\end{quote}

Its second character comes later; for now it is a word for the ear.

\begin{cousinweb}
\emph{yǒushì} is built from yǒu 'to have', the \zh{有} you write, and shì 'a matter'. Both are inherited from Middle Chinese, \emph{hjuwX} and \emph{dzriH} (Guangyun, in Baxter's transcription). The Guangyun glosses \emph{dzriH} 'to employ; to set up', and 'a matter' is the everyday sense, not the gloss.
\end{cousinweb}

\begin{grammarlens}[title={What you've built}]
\emph{yǒushì} is the polite reason that gives nothing away: something has come up. \emph{máng} says you are busy; \emph{yǒushì} says there is a particular thing.
\end{grammarlens}

\subsection*{Your turn}
\begin{itemize}
\raggedright
  \item \emph{Say it:} \emph{yǒushì}
  \item \emph{Say it:} \emph{yǒushì}, once more
  \item \emph{Say it:} \emph{wǒ jīntiān yǒushì} — I've got something on today.
  \item \emph{Recall:} say \emph{máng}, then \emph{yǒushì} again
\end{itemize}

\begin{tcolorbox}[breakable,colback=teal!4,colframe=teal!35!black,title={Before you move on}]
What does \emph{yǒushì} mean? (\textquotedblleft{}To have something on\textquotedblright{}.) Which written character does it open with? (\textbf{\zh{有}}.)
\end{tcolorbox}
\section[dǔchē]{dǔchē — a traffic jam}

\label{lesson:ZH-C95-duche}



\begin{quote}
Before the new one: last lesson's word, then two words for the road.

\begin{itemize}
\raggedright
  \item \emph{Recall:} say \emph{yǒushì}, to have something on — \textbf{R1}, one lesson back
  \item \emph{Recall:} say \zh{车} — \emph{chē}, a car
  \item \emph{Recall:} say \emph{lù}, a road — \textbf{R3}, durable
\end{itemize}
\end{quote}

\subsection*{You'll want to know: dǔchē}
\emph{dǔchē} — \textquotedblleft{}traffic jam; to be stuck in traffic\textquotedblright{}.

\begin{quote}
\emph{zǎoshang dǔchē} — There's a jam in the mornings.
\end{quote}

Its first character comes later; for now it is a word for the ear.

\begin{cousinweb}
\emph{dǔchē} is built from dǔ 'to block' and chē 'a vehicle', the \zh{车} you write. Both are inherited from Middle Chinese, \emph{tuX} and \emph{tsyhae} (Guangyun, in Baxter's transcription); the Guangyun glosses \emph{tuX} as 'a wall', and 'to block' is assumed from that, not quoted.
\end{cousinweb}

\begin{grammarlens}[title={What you've built}]
\emph{dǔchē} is a verb as well as a noun: \emph{zǎoshang dǔchē}, the traffic jams in the morning. No \zh{有} needed.
\end{grammarlens}

\subsection*{Your turn}
\begin{itemize}
\raggedright
  \item \emph{Say it:} \emph{dǔchē}
  \item \emph{Say it:} \emph{dǔchē}, once more
  \item \emph{Say it:} \emph{zǎoshang dǔchē} — There's a jam in the mornings.
  \item \emph{Recall:} say \emph{lù}, then \emph{dǔchē} again
\end{itemize}

\begin{tcolorbox}[breakable,colback=teal!4,colframe=teal!35!black,title={Before you move on}]
What does \emph{dǔchē} mean? (\textquotedblleft{}A traffic jam\textquotedblright{}.) Which written character does it end with? (\textbf{\zh{车}}.)
\end{tcolorbox}
\section[wàng]{wàng — to forget}

\label{lesson:ZH-C95-wang}



\begin{quote}
Before the new one: last lesson's word, then two words you might forget.

\begin{itemize}
\raggedright
  \item \emph{Recall:} say \emph{dǔchē}, a traffic jam — \textbf{R1}, one lesson back
  \item \emph{Recall:} say \emph{zuòyè}, homework — \textbf{R2}
  \item \emph{Recall:} say \zh{马上} — \emph{mǎshàng}, right away
\end{itemize}
\end{quote}

\subsection*{You'll want to know: wàng}
\emph{wàng} — \textquotedblleft{}to forget\textquotedblright{}.

\begin{quote}
\emph{wǒ wàng le} — I forgot.
\end{quote}

\begin{quote}
\emph{wǒ wàng le wǒ de yàoshi} — I forgot my key.
\end{quote}

Its character comes later; for now it is a word for the ear.

\begin{cousinweb}
\emph{wàng} is inherited from Middle Chinese \emph{mjangH} (Guangyun, in Baxter's transcription), glossed 'to leave behind; to forget'.
\end{cousinweb}

\begin{grammarlens}[title={What you've built}]
\emph{wàng le}: with \zh{了} after it, the forgetting has already happened, which is how it is almost always said.
\end{grammarlens}

\subsection*{Your turn}
\begin{itemize}
\raggedright
  \item \emph{Say it:} \emph{wàng}
  \item \emph{Say it:} \emph{wàng le}
  \item \emph{Say it:} \emph{wǒ wàng le} — I forgot.
  \item \emph{Say it:} \emph{wǒ wàng le wǒ de yàoshi} — I forgot my key.
  \item \emph{Recall:} say \emph{zuòyè}, then say you forgot it
\end{itemize}

\begin{tcolorbox}[breakable,colback=teal!4,colframe=teal!35!black,title={Before you move on}]
What does \emph{wàng} mean? (\textquotedblleft{}To forget\textquotedblright{}.) Say you forgot. (\emph{wǒ wàng le}.)
\end{tcolorbox}
\section[Practice]{\zh{你为什么不去}}

\label{lesson:ZH-C95-practice}



\begin{quote}
The chapter closes here, with one person asking why three times.

\begin{itemize}
\raggedright
  \item \emph{Recall:} say \emph{wàng}, to forget — \textbf{R1}, one lesson back
  \item \emph{Recall:} say \emph{láibují}, to run out of time — \textbf{R2}
  \item \emph{Recall:} say \emph{zuòyè}, homework — \textbf{R1}
\end{itemize}
\end{quote}

\subsection*{The exchange}
\begin{quote}
\textbf{A:} \emph{nǐ wèishénme chídào} — Why are you late?
\end{quote}

\begin{quote}
\textbf{B:} \emph{zǎoshang dǔchē} — The traffic this morning.
\end{quote}

\begin{quote}
\textbf{A:} \zh{你明天去吗}
\end{quote}

\begin{quote}
\emph{nǐ míngtiān qù ma} — Are you going tomorrow?
\end{quote}

\begin{quote}
\textbf{B:} \zh{我不去}
\end{quote}

\begin{quote}
\emph{wǒ bú qù} — I'm not going.
\end{quote}

\begin{quote}
\textbf{A:} \zh{为什么}
\end{quote}

\begin{quote}
\emph{wèishénme} — Why?
\end{quote}

\begin{quote}
\textbf{B:} \emph{wǒ yǒushì} — I've got something on.
\end{quote}

\begin{quote}
\textbf{A:} \emph{nǐ de zuòyè ne} — And your homework?
\end{quote}

\begin{quote}
\textbf{B:} \emph{wǒ wàng le} — I forgot.
\end{quote}

\zh{为什么} can stand alone as a whole question, as A's fifth line does. Each answer here is a reason with nothing in front of it; the next chapter gives the word that introduces one.

\begin{tcolorbox}[breakable,skin=enhanced,colback=orange!6,colframe=orange!45!black,boxrule=0.5pt,arc=1mm,left=6pt,right=6pt,top=4pt,bottom=4pt,fonttitle=\bfseries,title={Read this}]
\emph{Read it:} the exchange through once without stopping, then answer the question below
\end{tcolorbox}

\begin{grammarlens}[title={What you've built — asking why}]
\begin{itemize}
\raggedright
  \item \textbf{\zh{你为什么不去}} \emph{nǐ wèishénme bú qù} — why aren't you going?
  \item \textbf{\zh{为什么}} \emph{wèishénme} — why?
  \item \emph{wǒ yǒushì} — something has come up.
  \item \emph{zǎoshang dǔchē} — the traffic this morning.
\end{itemize}

\textbf{Listen.} \emph{Nǐ wèishénme bù chī fàn?}

\textbf{Speak.} Answer it with any reason you have, then ask somebody why.

\emph{Write it:} \zh{为} — the dot, the long fall, the hooked turn, the dot inside
\end{grammarlens}

\subsection*{Your turn}
\begin{itemize}
\raggedright
  \item \emph{Say it:} \zh{你为什么不去} — \emph{nǐ wèishénme bú qù} — Why aren't you going?
  \item \emph{Say it:} \emph{wǒ wàng le} — I forgot.
  \item \emph{Say it:} the whole exchange, both parts
  \item \emph{Write it:} \zh{为什么}
\end{itemize}

\begin{tcolorbox}[breakable,colback=teal!4,colframe=teal!35!black,title={Before you move on}]
Ask why somebody is not going. (\textbf{\zh{你为什么不去}}.) Say something has come up. (\emph{wǒ yǒushì}.) Blame the traffic. (\emph{zǎoshang dǔchē}.) Say you forgot. (\emph{wǒ wàng le}.)
\end{tcolorbox}
